% Source types: Mahbubur Rahman, Chapter 5 PDF pp.90--93.
\booksubheading{Recursive Function: ছোট সমস্যায় নিজেকে Call করা}

যে function তার কাজের মধ্যে নিজেকেই call করে, তাকে recursive function
বলে। দুটি বিষয় দরকার: কাজ থামানোর \textbf{base case} এবং আগের চেয়ে ছোট
সমস্যা নিয়ে পরের call। নতুন call-এ আলাদা local variable থাকে।

\textbf{উদাহরণ: Factorial-এর recursive রূপ।} Loop-ভিত্তিক সমাধান আগে
দেওয়া হয়েছে; এখানে function call ও return-এর ক্রম দেখানো হলো।
\(n\) একটি অঋণাত্মক পূর্ণসংখ্যা।
\begin{pythoncode}
def factorial(n):
    if n < 0:
        raise ValueError("n must be non-negative")
    if n == 0:
        return 1
    return n * factorial(n - 1)

print(factorial(4))
print(factorial(0))
\end{pythoncode}
\textbf{Output:} \py{24}, \py{1}। \py{raise ValueError(...)} এখানে
অগ্রহণযোগ্য input হলে ভুলের কারণ জানিয়ে ওই call থামায়।

\begin{tabularx}{\linewidth}{|Y|Y|}
\hline
\rowcolor{softBlue}\textbf{Call-এর সময়} & \textbf{ফিরে আসার সময়}\\\hline
\py{factorial(4)} অপেক্ষা করে \py{factorial(3)}-এর জন্য & \(4\times6=24\)\\\hline
\py{factorial(3)} অপেক্ষা করে \py{factorial(2)}-এর জন্য & \(3\times2=6\)\\\hline
\py{factorial(2)} অপেক্ষা করে \py{factorial(1)}-এর জন্য & \(2\times1=2\)\\\hline
\py{factorial(1)} অপেক্ষা করে \py{factorial(0)}-এর জন্য & \(1\times1=1\)\\\hline
\py{factorial(0)} base case-এ থামে & সরাসরি 1 return করে\\\hline
\end{tabularx}

Base case বাদ দিলে বা প্রতিবার একই \py{n} পাঠালে call থামার পথে এগোয় না।
প্রতিটি call-এর তথ্য memory-তে রাখতে হয়; তাই খুব বড় \py{n}-এর জন্য
Python-এর recursion limit-এ ভুল হতে পারে। এত বড় গণনায় আগের iterative
সমাধান উপযোগী। Recursion-এ আলাদা loop লেখা বাধ্যতামূলক নয়।

\booksubheading{এক File-এর Function অন্য Program-এ ব্যবহার}

পুনর্ব্যবহারের জন্য function একটি \py{.py} file-এ রাখা যায়। এই file-টিকে
module হিসেবে import করলে function আবার লিখতে হয় না। নিচের দুই file
একই folder-এ রাখো।

\textbf{প্রথম file: \py{geometry\_tools.py}}
\begin{pythoncode}
def rectangle_area(length, width):
    return length * width

def rectangle_perimeter(length, width):
    return 2 * (length + width)
\end{pythoncode}
\Needspace{10\baselineskip}
\textbf{দ্বিতীয় file: \py{use\_geometry.py}}
\begin{pythoncode}
import geometry_tools

print(geometry_tools.rectangle_area(7, 4))
print(geometry_tools.rectangle_perimeter(7, 4))
\end{pythoncode}
দ্বিতীয় file চালালে \textbf{Output:} \py{28}, \py{22}। Import করার সময়
module-এর নামের সঙ্গে \py{.py} লেখা হয় না। ক্ষেত্রফল ও পরিসীমার কাজ
একটি module-এ থাকায় অন্য প্রোগ্রাম থেকেও একই function ব্যবহার করা যায়।
\py{math.sqrt()}-এর মতোই module-এর নাম, dot ও function-এর নাম
দিয়ে call করা হয়েছে।

